\section{Thực trạng và định hướng quản lý chất lượng tại Trường Cao đẳng Kỹ Nghệ II}

Trong tiến trình đổi mới căn bản và toàn diện giáo dục nghề nghiệp tại Việt Nam, chất lượng không còn dừng lại ở một khẩu hiệu hành động hay một mục tiêu tĩnh mang tính hình thức, mà đã trở thành nhân tố quyết định sự tồn tại, uy tín thương hiệu và tính bền vững của mỗi cơ sở đào tạo. Sự ra đời của Thông tư số 14/2024/TT-BLĐTBXH ngày 31 tháng 12 năm 2024 quy định tiêu chí, tiêu chuẩn kiểm định chất lượng giáo dục nghề nghiệp đã đánh dấu bước chuyển mình quan trọng: từ phương thức tiếp cận kiểm tra mang tính cơ học, xử lý sự việc đã rồi sang cơ chế quản trị chất lượng dựa trên dữ liệu thực chất, nhấn mạnh trách nhiệm giải trình xã hội và sự hài lòng của các bên liên quan.

Đối với Trường Cao đẳng Kỹ Nghệ II tại Thành phố Hồ Chí Minh -- một trong những cơ sở giáo dục nghề nghiệp công lập tiên phong của cả nước thực hiện thí điểm cơ chế tự chủ tài chính toàn diện -- bài toán chất lượng mang tính sống còn. Tự chủ mang lại không gian sáng tạo và quyền tự quyết cao trong phân bổ nguồn lực, nhưng đồng thời cũng đặt nhà trường dưới áp lực cạnh tranh khốc liệt của thị trường lao động tại Vùng kinh tế trọng điểm phía Nam. Nhằm đáp ứng toàn diện bốn chuẩn đầu ra của học phần Quản lý chất lượng giáo dục nghề nghiệp (mã học phần ED6111), đề án tập trung nghiên cứu và giải quyết hệ thống bốn nội dung trọng tâm: phân tích thực trạng và định hướng cải tiến chất lượng cơ sở; thiết lập bộ tiêu chuẩn, công cụ đánh giá và thực hành tự đánh giá; thiết kế hệ thống đảm bảo chất lượng bên trong số hóa tích hợp tiêu chuẩn ISO 21001:2018; và xử lý xung đột đạo đức nghề nghiệp nhằm bảo đảm tính toàn vẹn của dữ liệu kiểm định.

\subsection{Bối cảnh phát triển và định vị chiến lược trong cơ chế tự chủ}

Trường Cao đẳng Kỹ Nghệ II, tiền thân là Trường Công nhân Kỹ thuật chế biến gỗ Trung ương được thành lập từ năm 1978 trực thuộc Bộ Lao động -- Thương binh và Xã hội, là một trong những cơ sở giáo dục nghề nghiệp công lập trọng điểm của khu vực phía Nam. Nhà trường tọa lạc tại vị trí cửa ngõ phía Đông của Thành phố Hồ Chí Minh, thuộc địa bàn Thành phố Thủ Đức. Vị trí địa lý này nằm tiếp giáp trực tiếp với mạng lưới các khu công nghiệp tập trung quy mô lớn như Khu Công nghệ cao Thành phố Hồ Chí Minh, Khu Chế xuất Linh Trung và Khu Công nghiệp Sóng Thần. Không gian liên kết kinh tế này tạo điều kiện thuận lợi để nhà trường xác lập mối quan hệ hợp tác đào tạo và tuyển dụng trực tiếp với hàng trăm doanh nghiệp sản xuất kỹ thuật hiện đại.

Về phương diện định vị chiến lược, Trường Cao đẳng Kỹ Nghệ II là một trong số ít các cơ sở giáo dục nghề nghiệp công lập trên toàn quốc được Thủ tướng Chính phủ phê duyệt thực hiện thí điểm cơ chế tự chủ tài chính toàn diện, bao gồm tự chủ toàn bộ chi thường xuyên và chi đầu tư theo tinh thần Nghị định số 60/2021/NĐ-CP. Sứ mạng của nhà trường được xác định nhất quán là đào tạo nguồn nhân lực kỹ thuật đa ngành có trình độ kỹ năng nghề xuất sắc, tác phong công nghiệp chuẩn mực và năng lực thích ứng linh hoạt với thị trường lao động trong nước cũng như quốc tế. Nhằm hiện thực hóa sứ mạng này, Chính sách chất lượng giai đoạn 2022 -- 2025 của nhà trường mang thông điệp chỉ đạo là đổi mới tư duy và tự chủ để phát triển bền vững.

Cơ chế tự chủ tài chính toàn diện vừa mang lại quyền tự quyết chủ động trong việc tổ chức bộ máy, tuyển dụng nhân sự và phát triển chương trình đào tạo, vừa đặt nhà trường trước áp lực sống còn về chỉ tiêu tuyển sinh và duy trì nguồn thu. Trong bối cảnh không còn nguồn ngân sách nhà nước cấp chi thường xuyên trực tiếp, thương hiệu của nhà trường phụ thuộc hoàn toàn vào mức độ hài lòng của người học và sự thừa nhận của các doanh nghiệp sử dụng lao động. Vì vậy, việc thiết lập một hệ thống quản trị chất lượng thực chất, tránh bệnh hình thức trên hồ sơ văn bản là đòi hỏi cấp bách và mang tính chiến lược của nhà trường.

\subsection{Môi trường vận hành và các yếu tố tác động đến chất lượng}

Để có cơ sở khoa học cho việc đề xuất các định hướng cải tiến chất lượng, việc xem xét thực trạng của Trường Cao đẳng Kỹ Nghệ II cần được đặt trong mối tương quan giữa các lực đẩy vĩ mô của môi trường kinh tế xã hội và các điều kiện nguồn lực nội tại của cơ sở giáo dục.

\subsubsection{Các yếu tố môi trường vĩ mô}

Môi trường thể chế và chính sách pháp luật hiện nay đang trải qua bước chuyển đổi căn bản với việc ban hành Thông tư số 14/2024/TT-BLĐTBXH quy định tiêu chí, tiêu chuẩn kiểm định chất lượng giáo dục nghề nghiệp. Văn bản quy phạm pháp luật này chuyển trọng tâm từ việc kiểm tra số lượng hồ sơ cơ học sang đánh giá chất lượng thực chất dựa trên chuẩn đầu ra và trách nhiệm giải trình xã hội của người đứng đầu cơ sở giáo dục. Về mặt kinh tế, Vùng kinh tế trọng điểm phía Nam đang dịch chuyển mạnh mẽ sang các ngành công nghiệp bán dẫn, tự động hóa, logistics thông minh và sản xuất bền vững. Xu hướng chuyển dịch kinh tế này đòi hỏi cơ sở giáo dục nghề nghiệp phải liên tục cập nhật công nghệ đào tạo, nếu không sinh viên tốt nghiệp sẽ đối mặt với nguy cơ bị đào thải khỏi thị trường lao động kỹ thuật cao.

Về phương diện xã hội, định kiến về bằng cấp đại học và tâm lý e ngại học nghề của phụ huynh vẫn là rào cản lớn đối với công tác phân luồng học sinh sau trung học cơ sở và trung học phổ thông. Đối tượng tuyển sinh ngày càng đa dạng, đòi hỏi các dịch vụ chăm sóc và tư vấn hỗ trợ người học phải được chuyên nghiệp hóa. Trên bình diện công nghệ, sự bùng nổ của trí tuệ nhân tạo, thiết bị kết nối vạn vật và công nghệ mô phỏng số đã làm thay đổi hoàn toàn phương thức giảng dạy nghề nghiệp truyền thống. Cuối cùng, cam kết quốc gia về phát thải ròng bằng không cùng các tiêu chuẩn lao động quốc tế trong các hiệp định thương mại tự do thế hệ mới đòi hỏi chuẩn đầu ra của nhà trường phải tích hợp kỹ năng xanh và năng lực ngoại ngữ chuyên ngành theo Khung trình độ quốc gia Việt Nam.

\begin{table}[htbp]
\centering
\small
\setstretch{1.2}
\caption{\textbf{Ma trận phân tích các yếu tố vĩ mô tác động đến Trường Cao đẳng Kỹ Nghệ II}}
\vspace{0.2cm}
\begin{tabularx}{\textwidth}{L{2.8cm} X X}
\toprule
\textbf{Yếu tố vĩ mô} & \textbf{Nội dung tác động đến cơ sở giáo dục} & \textbf{Hàm ý đối với quản lý chất lượng} \\
\midrule
Chính trị và Pháp lý & Khung khổ pháp lý mới với Thông tư số 14/2024/TT-BLĐTBXH đặt trọng tâm vào chuẩn đầu ra và trách nhiệm giải trình. & Yêu cầu chuẩn hóa toàn diện kho minh chứng, chấm dứt việc lập hồ sơ đối phó và gắn trách nhiệm pháp lý với cam kết chất lượng. \\
\addlinespace
Kinh tế & Vùng Đông Nam Bộ chuyển dịch mạnh sang tự động hóa, bán dẫn và logistics thông minh. & Chương trình đào tạo và thiết bị xưởng phải được nâng cấp liên tục để bảo đảm tỷ lệ sinh viên có việc làm đúng chuyên môn. \\
\addlinespace
Xã hội & Tâm lý chuộng bằng cấp đại học vẫn phổ biến; sự phân hóa về hoàn cảnh gia đình của người học gia tăng. & Công tác tư vấn học đường, hỗ trợ sinh viên khó khăn và quản lý duy trì sĩ số trở thành yếu tố quyết định hiệu quả giáo dục. \\
\addlinespace
Công nghệ & Trí tuệ nhân tạo, tự động hóa và công nghệ số hóa tái định hình các thao tác nghề nghiệp tại nhà xưởng. & Bắt buộc phải tích hợp chuẩn đầu ra năng lực số vào các ngành nghề và số hóa toàn diện quy trình kiểm tra đánh giá. \\
\addlinespace
Môi trường và Hội nhập & Chiến lược tăng trưởng xanh và các cam kết kỹ năng nghề tiêu chuẩn khu vực Đông Nam Á. & Đưa tiêu chuẩn an toàn sinh thái vào giảng dạy và đối sánh chuẩn đầu ra theo Khung trình độ quốc gia Việt Nam. \\
\bottomrule
\end{tabularx}
\end{table}

\subsubsection{Nguồn lực nội tại và vị thế cạnh tranh}

Xét trên phương diện điểm mạnh nội tại, Trường Cao đẳng Kỹ Nghệ II sở hữu ưu thế vượt trội về vị trí địa lý đắc địa tiếp giáp Khu Công nghệ cao Thành phố Hồ Chí Minh, giúp tạo lập mạng lưới doanh nghiệp đối tác tiếp nhận thực tập rộng lớn. Cơ sở vật chất của nhà trường được đầu tư đồng bộ với nhiều nhà xưởng thực hành hiện đại đáp ứng các tiêu chuẩn công nghiệp. Đội ngũ giảng viên có kỹ năng nghề thực tế vững vàng, hầu hết đều đạt chuẩn chức danh nghề nghiệp và có tỷ lệ tham gia bồi dưỡng tại doanh nghiệp hàng năm cao. Đặc biệt, nhà trường đã ban hành được hệ thống gồm sáu mươi lăm quy trình vận hành và kiểm soát nội bộ, tạo nền tảng văn bản hóa vững chắc cho toàn bộ hoạt động quản trị.

Tuy nhiên, phân tích điểm yếu nội tại cho thấy hệ thống dữ liệu phục vụ quản lý chất lượng của nhà trường vẫn trong tình trạng phân mảnh nghiêm trọng. Việc lưu trữ hồ sơ giấy tờ riêng rẽ tại các khoa và phòng chức năng làm cản trở quá trình đối soát dữ liệu liên thông theo thời gian thực. Cơ chế khảo sát ý kiến người học và nhà tuyển dụng vẫn chủ yếu thực hiện vào cuối học kỳ hoặc sau khi tốt nghiệp, gây ra độ trễ thông tin lớn và làm mất đi cơ hội điều chỉnh phương pháp giảng dạy ngay trong quá trình đào tạo. Thêm vào đó, nhận thức về công tác đảm bảo chất lượng tại một số khoa chuyên môn vẫn mang tính thụ động, xem đây là nhiệm vụ riêng của Phòng Đảm bảo chất lượng.

Về cơ hội phát triển, cơ chế tự chủ tài chính toàn diện cho phép nhà trường linh hoạt mở các ngành đào tạo mới đón đầu xu hướng công nghệ số và liên kết đào tạo song hành cùng các tập đoàn đa quốc gia. Quy định kiểm định mới theo Thông tư số 14/2024/TT-BLĐTBXH cũng tạo ra hành lang pháp lý minh bạch để khẳng định uy tín thương hiệu của những cơ sở có chất lượng thực chất. Ngược lại, thách thức bên ngoài lớn nhất xuất phát từ sự cạnh tranh tuyển sinh gay gắt của các trường đại học thông qua phương thức xét tuyển học bạ mở rộng, cùng với áp lực tài chính rất lớn trong việc tạo nguồn vốn tự tích lũy để bảo trì và thay mới các thiết bị công nghệ cao có chu kỳ hao mòn nhanh.

\subsection{Thực trạng công tác đảm bảo chất lượng qua kết quả tự đánh giá}

Công tác tự đánh giá chất lượng cơ sở giáo dục nghề nghiệp của Trường Cao đẳng Kỹ Nghệ II được thực hiện vào cuối năm 2024 và chính thức phê duyệt tại Báo cáo số 242/BC-CĐKNII ngày 31 tháng 01 năm 2025 của Hiệu trưởng, áp dụng bộ tiêu chí ban hành kèm theo Thông tư số 14/2024/TT-BLĐTBXH. Kết quả tự đánh giá ghi nhận nhà trường đạt điểm số tuyệt đối ở toàn bộ các tiêu chí và tiêu chuẩn kiểm định.

\begin{table}[htbp]
\centering
\small
\setstretch{1.2}
\caption{\textbf{Tổng hợp kết quả tự đánh giá chất lượng cơ sở năm 2025 của Trường Cao đẳng Kỹ Nghệ II}}
\vspace{0.2cm}
\begin{tabularx}{\textwidth}{C{1.0cm} X C{2.2cm} C{1.8cm} C{2.2cm} C{1.8cm}}
\toprule
\textbf{STT} & \textbf{Tên tiêu chí kiểm định chất lượng} & \textbf{Số lượng tiêu chuẩn} & \textbf{Điểm tối đa} & \textbf{Điểm tự đánh giá} & \textbf{Tỷ lệ đạt (\%)} \\
\midrule
1 & Sứ mạng, mục tiêu, tổ chức và quản lý & 5 & 10 & 10 & 100\% \\
2 & Hoạt động đào tạo & 8 & 16 & 16 & 100\% \\
3 & Nhà giáo, cán bộ quản lý, viên chức và người lao động & 7 & 14 & 14 & 100\% \\
4 & Chương trình đào tạo và giáo trình & 7 & 14 & 14 & 100\% \\
5 & Cơ sở vật chất, thiết bị đào tạo và thư viện & 8 & 16 & 16 & 100\% \\
6 & Nghiên cứu khoa học, chuyển giao công nghệ và hợp tác quốc tế & 4 & 8 & 8 & 100\% \\
7 & Dịch vụ người học & 5 & 10 & 10 & 100\% \\
8 & Giám sát, đánh giá chất lượng & 6 & 12 & 12 & 100\% \\
\midrule
\multicolumn{2}{l}{\textbf{Tổng hệ thống đánh giá cơ sở}} & \textbf{50} & \textbf{100} & \textbf{100} & \textbf{100\%} \\
\bottomrule
\end{tabularx}
\end{table}

Xét trên phương diện tuân thủ pháp luật, kết quả tự đánh giá một trăm điểm khẳng định năng lực chuẩn bị hồ sơ minh chứng rất bài bản của tập thể nhà trường. Các tiêu chuẩn về tổ chức quản lý, trình độ đội ngũ nhà giáo, tiến độ rà soát chương trình đào tạo và diện tích phòng học, nhà xưởng bình quân trên mỗi người học đều được minh chứng đầy đủ bằng các quyết định hành chính, văn bằng chứng chỉ và sổ sách kỹ thuật hợp lệ. Chính sách miễn giảm học phí và cấp phát học bổng khuyến khích học tập cho người học được thực hiện công khai, minh bạch thông qua tài khoản ngân hàng, không để xảy ra khiếu nại phát sinh.

\subsection{Các khoảng trống chất lượng cốt lõi và rủi ro quản trị}

Dưới lăng kính của Total Quality Management (Quản lý chất lượng toàn diện -- TQM) và nguyên tắc quản trị rủi ro tại Điều khoản 6.1 của tiêu chuẩn quốc tế ISO 21001:2018, việc đạt điểm chuẩn tuyệt đối trên hồ sơ tự đánh giá mới chỉ phản ánh mức độ tuân thủ quy chuẩn hành chính, chưa bảo đảm việc loại trừ các rủi ro chất lượng trong vận hành thực tế. Khi đối chiếu Báo cáo tự đánh giá với Bảng mười tám Mục tiêu chất lượng năm 2025 do Hiệu trưởng ban hành, ba khoảng trống chất lượng bản chất đã bộc lộ rõ rệt.

Khoảng trống thứ nhất nằm ở tư duy quản trị chấp nhận tỷ lệ hao hụt người học ở mức cao và sự thiếu vắng cơ chế cảnh báo sớm. Cụ thể, Mục tiêu chất lượng số 3 năm 2025 của nhà trường chỉ đặt chỉ tiêu duy trì sĩ số học sinh, sinh viên từ bảy mươi lăm phần trăm trở lên. Điều này đồng nghĩa với việc hệ thống quản trị của trường mặc nhiên chấp nhận tỷ lệ bỏ học, thôi học có thể lên tới hai mươi lăm phần trăm trong mỗi khóa đào tạo. Dưới góc độ kinh tế giáo dục và trách nhiệm xã hội, việc để một phần tư người học rời bỏ nhà trường giữa chừng gây thất thoát tài chính lớn cho trường tự chủ và lãng phí nguồn lực đầu tư của xã hội. Nguyên nhân sâu xa là nhà trường chỉ ghi nhận số liệu sau khi sự việc đã hoàn thành bằng các quyết định buộc thôi học, hoàn toàn thiếu một hệ thống cảnh báo sớm có khả năng tự động nhận diện các sinh viên có nguy cơ bỏ học thông qua dữ liệu chuyên cần và kết quả kiểm tra định kỳ để can thiệp hỗ trợ kịp thời.

Khoảng trống thứ hai liên quan đến độ trễ thông tin trong các quy trình thu thập phản hồi của các bên liên quan. Mặc dù các quy trình khảo sát người học có mã hiệu QT/BĐCL/02, quy trình khảo sát giảng viên có mã hiệu QT/BĐCL/03 và quy trình khảo sát việc làm có mã hiệu ĐT/QT25/KSVL đã được ban hành đầy đủ, phương thức thực hiện trên thực tế vẫn chủ yếu dùng bảng hỏi định kỳ phát vào cuối học kỳ hoặc sau khi người học đã tốt nghiệp một năm. Cách tiếp cận này mang bản chất của kiểm tra chất lượng hồi cứu nhằm tổng kết quá khứ hơn là đảm bảo chất lượng chủ động. Giảng viên không thể điều chỉnh phương pháp sư phạm cho chính lớp học đang đảm nhận khi chỉ nhận được kết quả phản hồi sau khi môn học đã kết thúc. Nhà trường hoàn toàn thiếu các công cụ khảo sát tức thời sau từng bài học để điều chỉnh chất lượng đào tạo liên tục.

Khoảng trống thứ ba xuất phát từ mức độ phân mảnh trong hệ thống dữ liệu quản lý và cơ chế kiểm chứng chéo minh chứng. Mặc dù nhà trường đã xây dựng danh mục sáu mươi lăm quy trình vận hành bằng văn bản, các nguồn dữ liệu tác nghiệp thực tế vẫn bị cô lập trong các phần mềm và sổ sách riêng lẻ giữa Phòng Đào tạo, Phòng Công tác sinh viên, Phòng Kế toán và các khoa chuyên môn. Tình trạng này khiến công tác trích xuất và kiểm tra tính xác thực của dữ liệu tốn nhiều công sức, đồng thời tiềm ẩn rủi ro phát sinh khoảng cách sai lệch giữa số liệu báo cáo minh chứng trên giấy tờ với mức độ thụ hưởng chất lượng thực tế của sinh viên tại các xưởng máy và giảng đường.

\subsection{Định hướng chiến lược phát triển chất lượng bền vững}

Căn cứ vào kết quả phân tích các yếu tố môi trường và ba khoảng trống chất lượng cốt lõi, bốn định hướng chiến lược cải tiến chất lượng mang tính đột phá được xác lập cho Trường Cao đẳng Kỹ Nghệ II trong giai đoạn 2026 -- 2030.

Định hướng thứ nhất là chuyển đổi căn bản mô hình quản trị chất lượng từ phương thức đối phó tuân thủ sang xây dựng hệ thống đảm bảo chất lượng bên trong được thúc đẩy bởi công nghệ số, tích hợp chặt chẽ triết lý Quản lý chất lượng toàn diện và tiêu chuẩn quốc tế ISO 21001:2018. Nhà trường cần xây dựng trung tâm điều hành dữ liệu chất lượng liên thông toàn trường, số hóa toàn bộ dữ liệu tác nghiệp của sáu mươi lăm quy trình nội bộ lên nền tảng đám mây nhằm bảo đảm mọi minh chứng phục vụ kiểm định đều có dấu vết số hóa khách quan và không thể bị can thiệp làm sai lệch.

Định hướng thứ hai là thiết lập cơ chế kiểm soát chủ động biến động người học, phấn đấu nâng chỉ tiêu duy trì sĩ số toàn khóa lên mức tối thiểu tám mươi tám phần trăm, tương ứng với việc khống chế tỷ lệ bỏ học dưới mười hai phần trăm. Nhà trường cần xóa bỏ tư duy chấp nhận tỷ lệ hao hụt hai mươi lăm phần trăm bằng việc đưa vào vận hành hệ thống cảnh báo sớm tự động trên phần mềm quản trị đào tạo. Khi sinh viên có dấu hiệu bất thường về chuyên cần hoặc kết quả học tập sa sút, hệ thống sẽ tự động kích hoạt mạng lưới cố vấn học tập và bộ phận hỗ trợ sinh viên để can thiệp, tư vấn tâm lý và hỗ trợ tài chính trước khi người học rơi vào tình trạng buông xuôi.

Định hướng thứ ba là rút ngắn chu kỳ phản hồi chất lượng thông qua việc ứng dụng các công cụ khảo sát tức thời trên nền tảng ứng dụng di động của nhà trường. Thay thế các bảng hỏi dài dòng cuối kỳ bằng các phiếu khảo sát ngắn gọn từ ba đến năm câu hỏi được gửi tự động sau các bài thực hành trọng điểm, giúp giảng viên và cán bộ quản lý xưởng tiếp nhận phản hồi trong vòng hai mươi tư giờ để khắc phục ngay các bất cập sư phạm hoặc sự cố kỹ thuật. Song song với đó, nhà trường cần xây dựng văn hóa chất lượng tự giác, khuyến khích cán bộ, giảng viên và người học chủ động phát hiện và báo cáo các điểm nghẽn quy trình trên tinh thần cải tiến liên tục không đổ lỗi.

Định hướng thứ tư là thực hiện rà soát và cấu trúc lại toàn bộ hệ thống chuẩn đầu ra các ngành nghề đào tạo, tích hợp bắt buộc chuẩn đầu ra về năng lực số và kỹ năng chuyển đổi xanh. Nhà trường cần phối hợp chặt chẽ với các doanh nghiệp tại Khu Công nghệ cao Thành phố Hồ Chí Minh để đối sánh chương trình với tiêu chuẩn nghề nghiệp tiên tiến, đưa các học phần về tự động hóa, điều khiển số, an toàn dữ liệu và bảo vệ môi trường vào cấu trúc đào tạo bắt buộc, tạo lợi thế cạnh tranh vượt trội cho sinh viên sau khi tốt nghiệp.
